% ================
\section{Methodology}\label{sec:method}

% We model semantic consensus via a weighted Fréchet mean of answer embeddings, where the weighting distribution can optionally incorporate additional signals such as frequency or model confidence. The resulting semantic center is then used to rank candidates according to their distances to it. We provide the detailed pseudo-code for our method in Appendix Algorithm~\ref{alg:rds}.

\subsection{Problem Statement}\label{sec:problem}

Given a prompt $x$, we sample $N$ independent generations $\{y_1, \dots, y_N\}$ from a language model using multinomial decoding. Each generation is mapped to a final answer $\{a_i\}_{i=1}^N$, optionally associated with generation likelihoods $\{p(y_i | x)\}_{i=1}^N$.
Our goal is to identify the most reliable answer among the candidates. Formally, we consider the problem:
\begin{align}
    a^* = \arg\max_{a_i} \; Q(a_i|x),
\end{align}
where $Q(a_i|x)$ denote an unknown quality function reflecting the correctness or reliability of an answer, which must be estimated from a set of sampled generations. 
% Existing approaches, such as self-consistency~\citep{wang2022self}, approximate $Q(a_i \mid x)$ using answer frequency, implicitly assuming that correct answers appear more often. However, such discrete aggregation ignores semantic relationships between answers and treats all disagreements equally, motivating the need for a more expressive estimator of answer quality.

\subsection{Answer Selection via Geometric Consensus}~\label{sec:center}

To approximate the latent quality function $Q(a_i|x)$, we propose to model agreement among answers in a continuous semantic space. Formally, we embed each final answer into a vector space:
\begin{align}
    \mathbf{u}_i = \mathbf{E}(a_i),
\end{align}
where $\mathbf{E}$ is a sentence embedding model. 
Intuitively, in the context of best-of-$N$, sampled answers form a distribution of semantic interpretations of the same query, and a representative solution should reflect their overall agreement. The Fréchet mean~\citep{fletcher} defines such a centroid as the point that minimizes the expected squared distance to all candidate embeddings in the semantic space. Given embeddings $\{\mathbf{u}_i\}_{i=1}^N \subset \mathbb{R}^d$ with a probability distribution $P{=}\{p_i\}_{i=1}^N$, the semantic center is defined as the Fréchet mean under squared Euclidean loss:
\begin{align}
    \mathbf{c}(P) = \arg\min_{\mathbf{z}} \sum_{i=1}^N p_i \|\mathbf{u}_i - \mathbf{z}\|_2^2,
        \label{eq:center}
\end{align}
This admits the closed-form solution:
\begin{align}
    \mathbf{c}(P) = \sum_{i=1}^N p_i \mathbf{u}_i,
\end{align}
(see Appendix~\ref{app:proof_center} for details).

We adopt the squared Euclidean loss for its simple closed-form solution. When the center is restricted to the discrete set $\{\mathbf{u}_i\}_{i=1}^N$, the problem reduces to selecting a representative sample:
\begin{align}
\mathbf{c}_{\text{medoid}}(P)
= \arg\min_{\mathbf{u}_j} \sum_{i=1}^N p_i \|\mathbf{u}_i - \mathbf{u}_j\|_2^2,
\end{align}
which corresponds to a weighted medoid. This can be viewed as a discrete counterpart of the Fréchet mean, where the solution is constrained to be one of the observed samples.

Different choices of $P$ (Table~\ref{tab:rds_variants}) induce different notions of semantic consensus. In particular, the uniform distribution recovers the standard unweighted center, while frequency-based or likelihood-based distributions incorporate empirical or model-derived reliability signals.
Conceptually, $\mathbf{c}(P)$ represents a consensus embedding of the generated answers, with semantically similar answers pulling the center toward their shared meaning and divergent or noisy samples contributing proportionally to their weights in $P$ (see Figure~\ref{fig:overview}).
% Intuitively, $\mathbf{c}(P)$ represents a consensus embedding of the generated answers. Semantically similar answers pull the center toward their shared meaning, while divergent or noisy samples exert influence proportional to their weights in $P$ (see Figure~\ref{fig:overview}).

\subsection{Ranking via Radial Consensus Score}\label{sec:main-rcs}

Given a semantic center $\mathbf{c}(P)$, we define a ranking function based on the radial dispersion of each candidate embedding relative to this center.
For each candidate embedding $\mathbf{u}_i$, its \textbf{Radial Consensus Score} (RCS) is defined as:
\begin{align}
    \mathrm{RCS}_i(P) = \|\mathbf{u}_i - \mathbf{c}(P)\|_2.
\end{align}
% where $\mathbf{c}(P)$ is defined in Section~\ref{sec:center}.
RCS measures how semantically aligned each candidate is with the global consensus induced by the distribution $P$: smaller values indicate stronger agreement with the collective semantic structure, while larger values suggest divergence or potential noise. 
We adopt the $\ell_2$ norm for its simplicity and consistency with the center formulation (Equation~\ref{eq:center}), as the choice of norm does not materially affect the relative ranking of candidates in practice.

\begin{table*}[t]
\centering
\caption{Different instantiations of the distribution $P$ in RCS, where $f_i$ and $\Pr(y_i|\mathbf{x})$ denote the empirical frequency and model probability of candidate $y_i$, respectively. For the frequency-weighted variant, the index $i$ runs over the set of \emph{unique} answers, so that $\sum_j f_j = N$. For the cosine-similarity-weighted variant, $\cos(\mathbf{u}_i,\mathbf{u}_j)$ denotes the cosine similarity between $L_2$-normalized embeddings.}
\label{tab:rds_variants}
\begin{tabular}{l l l}
\toprule
\textbf{Variant} & \textbf{Distribution $P$} & \textbf{Note} \\
\midrule
Uniform & $p_i = \frac{1}{N}$ & Black-box\\
Frequency-weighted & $p_i = \frac{f_i}{\sum_j f_j}$ & Black-box\\
Probability-weighted & $p_i \propto \Pr(y_i|\mathbf{x})$ & White-box (calibrated)\\
Cosine-weighted & $p_i \propto \frac{1}{N}\sum_{j=1}^N \cos(\mathbf{u}_i,\mathbf{u}_j)$ & Black-box\\
\bottomrule
\end{tabular}
\end{table*}

The final prediction is selected as the candidate with minimum radial dispersion to the estimated consensus:
\begin{align}
    a^*(P) = \arg\min_{a_i} \mathrm{RCS}_i(P).
\end{align}

In discrete settings, the above objective admits an equivalent \textbf{weighted medoid} formulation:
\begin{align}
    a^*(P) = \arg\min_{a_i} \sum_{j} P(a_j)\,\|\mathbf{u}_i - \mathbf{u}_j\|_2^2.
\end{align}
This selects the candidate whose embedding minimizes the expected \emph{squared} distance to others under $P$, consistent with the squared-Euclidean Fréchet mean objective (Equation~\ref{eq:center}). When $P$ is uniform, this reduces to the standard medoid.


\paragraph{Theoretical Analysis.}
To formally characterize the behavior of RCS, we model the answer embeddings as draws from a Gaussian noise distribution centered at the true semantic center--that is, $\mathbf{u}_i = \mathbf{c}^* + \boldsymbol{\varepsilon}_i$ with $\boldsymbol{\varepsilon}_i \sim \mathcal{N}(\mathbf{0}, \sigma^2\mathbf{I})$. 
Under this generative assumption, we establish two propositions showing that RCS is strictly superior to majority voting, formulated as follows. The uniform, frequency, and probability-weighted variants admit direct proofs via the strong law of large numbers; the cosine-similarity-weighted variant satisfies the same asymptotic guarantee under an additional regularity condition (Appendix~\ref{app:proof_prop1}).
% with detailed proofs deferred to Appendix~\ref{app:proof_prop1} and~\ref{app:proof_prop2}.

\begin{prop}[\textbf{Asymptotic Superiority}]
Under the Gaussian generative model, for \emph{any} of the four distributions \(P\) (Table~\ref{tab:rds_variants}), as \(N \to \infty\),
\begin{align}
\mathbf{c}(P) \xrightarrow{a.s.} \mathbf{c}^*,
\end{align}
by the strong law of large numbers. Consequently,
\begin{align}
\mathbb{P}(\text{RCS selects $a^*$}) \to 1, \quad \mathbb{P}(\text{voting selects $a^*$}) \leq \frac{1}{M},
\end{align}
where $a^*$ is the true best answer (corresponding to $\mathbf{c}^*$) and \(M \geq 2\) is the number of distinct semantic clusters (typically \(M \gg 1\) in open-ended generation). Thus RCS is asymptotically strictly superior to majority voting whenever lexical variance (due to paraphrasing) exceeds semantic variance.
% \qed
\end{prop}
\begin{proof}
    See Appendix~\ref{app:proof_prop1}.
\end{proof}
% \qed

\begin{prop}[\textbf{Finite-\(N\) Gap with 2-Cluster Toy Model}]
Consider the following two-cluster setting: (1) Cluster A (true high-quality): \(k \geq N/2\) samples centered at \(\mathbf{c}^* + \boldsymbol{\delta}\), \(\|\boldsymbol{\delta}\|_2\) small; and (2) Cluster B (spurious): \(N-k\) samples centered at \(\mathbf{c}^* + \boldsymbol{\Delta}\), \(\|\boldsymbol{\Delta}\|_2 \gg \|\boldsymbol{\delta}\|_2\).

Under the frequency-weighted distribution, the empirical center \(\mathbf{c}(P)\) is geometrically biased toward Cluster A when \(k \geq N/2\). Approximating each cluster's candidates at their expected locations, RCS selects Cluster A with probability
\begin{align}
\mathbb{P}(\text{RCS correct})
= \Phi\left( \frac{(2k-N)\bigl(\|\boldsymbol{\Delta}\|_2 - \|\boldsymbol{\delta}\|_2\bigr)}{2\sigma\sqrt{N}} \right),
\end{align}
where \(\Phi\) is the standard normal CDF. Under exact string matching with \(m \geq 2\) distinct surface forms per semantic cluster, the most frequent correct form receives at most \(k/m\) out of \(N\) votes; a single random sample is correct with probability \(k/N\).

Therefore,
\begin{align}
\mathbb{P}(\text{RCS correct}) > \frac{k}{N}
\end{align}
whenever the semantic gap satisfies
\begin{align}
\|\boldsymbol{\Delta}\|_2 - \|\boldsymbol{\delta}\|_2 > C\frac{\sigma}{\sqrt{N}}, \quad C > 0.
\end{align}
This property holds for the uniform and probability-weighted variants.
\begin{remark}
When \(k < N/2\), the center \(\mathbf{c}(P)\) is biased toward Cluster B and the sufficient condition above does not hold. Proposition 1 still guarantees \(\mathbf{c}(P) \xrightarrow{a.s.} \mathbf{c}^*\) as \(N \to \infty\), ensuring asymptotic correctness for all \(k\).
\end{remark}
\end{prop}
\begin{proof}
    See Appendix~\ref{app:proof_prop2}.
\end{proof}

Overall, the choice of $P$ controls how consensus is formed: uniform weighting reflects pure geometric agreement, frequency/probability weighting emphasizes stable or high-confidence responses, and cosine-similarity weighting uses mutual embedding agreement---all naturally suppressing inconsistent outputs. RCS$_\text{uni}$, RCS$_\text{freq}$, and RCS$_\text{cosine}$ operate purely on model outputs and embeddings, making them applicable in black-box settings.
While RCS is defined over candidate-level embeddings, one may alternatively consider aggregating representations over full generation trajectories. However, as shown in Section~\ref{sec:embedding-collapse}, such approaches suffer from representation collapse and lead to degraded performance, highlighting the importance of focusing on final-answer semantics. 
We provide the detailed pseudo-code for our method in Appendix Algorithm~\ref{alg:rds}.